\section
[لم مازارت]
{عدد پوششی}
سه بخش این سوال ارتباطی با یکدیگر ندارند.
\subsection{عدد پوششی خارجی}
فرض کنید مجموعه‌ی 
$A$
و متر $d$ روی اعضای $A$ داده‌شده باشد. مجموعه‌ی 
$K$
را یک زیرمجموعه‌ از
$A$
در نظر بگیرید. در کلاس درس، با مفهوم عدد پوششی آشنا شدیم.

\begin{itemize}
\item
در یک تعریف عدد پوششی، فرض بر این است که مراکز گوی‌های پوشش‌دهنده‌ی مجموعه‌ی 
$K$،
 یعنی نقاط
$x_i$،
همگی اعضای مجموعه‌ی 
$K$
باشند. عدد پوششی به دست آمده از این تعریف را
$\mathcal{N}(K, d, \epsilon)$
می‌نامیم که 
$\epsilon$
شعاع گوی‌های پوشش‌دهنده است.
\item
در تعریف دوم، مراکز گوی‌های پوشش دهنده، الزامی به عضو $K$ بودن ندارند و می‌توانند هر عضو دلخواه $A$ باشند. عدد پوششی به دست آمده از این تعریف را 
$\mathcal{N}^{ext}(K, d, \epsilon)$
می‌نامیم.
\end{itemize}

به توجه به این دو تعریف، نشان دهید:
$$\mathcal{N}^{ext}(K, d, \epsilon) \leq \mathcal{N}(K, d, \epsilon) \leq \mathcal{N}^{ext}(K, d, \epsilon/2)$$

\subsection{یکنوایی عدد پوششی}
در ابتدا، مثال نقضی برای نامساوی زیر بزنید
$$L \subset K \implies \mathcal{N}(L, d, \epsilon) \leq \mathcal{N}(K, d, \epsilon)$$
سپس تلاش نمایید نامساوی زیر را ثابت کنید
$$L \subset K \implies \mathcal{N}(L, d, \epsilon) \leq \mathcal{N}(K, d, \epsilon/2).$$

\subsection
[مکعب همینگ]
{پوشش و گنجایش مکعب همینگ}
مجموعه‌ی
$K = \{0, 1\}^n$
را با فاصله‌ی همینگ در نظر بگیرید.  نشان دهید که برای هر عدد حسابی
$m \leq n$
داریم:
$$\frac{2^n}{\sum_{k = 1}^{m} {n \choose k}} \leq \mathcal{N}(K, d_H, m) \leq \mathcal{M}(K, d_H, m) \leq \frac{2^n}{\sum_{k = 0}^{[\frac{m}{2}]} {n \choose k}}$$
که در آن $d_H$ متر همینگ، 
$\mathcal{N}$
عدد پوششی و 
$\mathcal{M}$
عدد گنجایشی است.
